% Cahier des charges — Tournoi d'échecs en direct
% Compiler : pdfLaTeX
\documentclass[11pt,a4paper]{article}

\usepackage[T1]{fontenc}
\usepackage[utf8]{inputenc}
\usepackage[french]{babel}
\usepackage{helvet}
\renewcommand{\familydefault}{\sfdefault}

\usepackage[
  a4paper,
  top=2.15cm,
  bottom=2.15cm,
  left=2.15cm,
  right=2.15cm,
  headheight=22pt
]{geometry}

\usepackage{xcolor}
\usepackage{booktabs}
\usepackage{tabularx}
\usepackage{array}
\usepackage{enumitem}
\usepackage{titlesec}
\usepackage{fancyhdr}
\usepackage{microtype}
\usepackage{hyperref}

\definecolor{ink}{HTML}{1C1C1C}
\definecolor{muted}{HTML}{6A6A6A}
\definecolor{line}{HTML}{E6E6E6}
\definecolor{soft}{HTML}{F4F4F4}
\definecolor{rule}{HTML}{111111}

\hypersetup{
  colorlinks=true,
  linkcolor=ink,
  urlcolor=ink,
  pdftitle={Cahier des charges — Tournoi d'échecs en direct},
  pdfauthor={Raouf Abdallah}
}

\pagestyle{fancy}
\fancyhf{}
\fancyhead[L]{\small\color{muted}Tournoi d'échecs en direct}
\fancyhead[R]{\small\color{muted}Cahier des charges}
\fancyfoot[C]{\small\color{muted}\thepage}
\renewcommand{\headrulewidth}{0.4pt}
\renewcommand{\headrule}{%
  \hbox to\headwidth{\color{line}\leaders\hrule height \headrulewidth\hfill}%
}
\renewcommand{\footrulewidth}{0pt}

\fancypagestyle{opening}{%
  \fancyhf{}%
  \fancyfoot[C]{\small\color{muted}\thepage}%
  \renewcommand{\headrulewidth}{0pt}%
}

\setlist[itemize]{leftmargin=1.15em,itemsep=2.5pt,topsep=4pt,parsep=0pt,partopsep=0pt}
\setlist[enumerate]{leftmargin=1.35em,itemsep=3pt,topsep=4pt,parsep=0pt,partopsep=0pt}

\titleformat{\section}
  {\large\bfseries\color{ink}}
  {\thesection}
  {0.75em}
  {}
  [\vspace{2pt}{\color{line}\titlerule[0.4pt]}]
\titleformat{\subsection}
  {\normalsize\bfseries\color{ink}}
  {\thesubsection}
  {0.6em}
  {}
\titlespacing*{\section}{0pt}{16pt}{8pt}
\titlespacing*{\subsection}{0pt}{11pt}{4pt}

\setcounter{tocdepth}{1}
\setcounter{secnumdepth}{2}

\newcommand{\panel}[2]{%
  \par\vspace{6pt}%
  {\setlength{\fboxsep}{9pt}%
   \noindent\colorbox{soft}{%
     \parbox{\dimexpr\linewidth-2\fboxsep\relax}{#1\small #2}%
   }}%
  \par\vspace{6pt}%
}
\newcommand{\softbox}[1]{\panel{\ttfamily}{#1}}
\newcommand{\quotebox}[1]{\panel{}{#1}}

\newcommand{\field}[1]{{\ttfamily #1}}

\newcolumntype{L}{>{\raggedright\arraybackslash}X}
\newcolumntype{Y}{>{\raggedright\arraybackslash}p{4.15cm}}

\setlength{\parindent}{0pt}
\setlength{\parskip}{0.45em}
\raggedbottom
\color{ink}

\begin{document}
\thispagestyle{opening}

{\small\color{muted}Cahier des charges \textperiodcentered\ octobre 2026}

\vspace{8pt}
{\fontsize{24}{28}\selectfont\bfseries Tournoi d'échecs en direct\par}

\vspace{6pt}
{\large\color{muted}
Partie en direct, spectateur et rejeu.\\
Java, Spring Boot, PostgreSQL, Kafka, Angular.\par}

\vspace{10pt}
{\color{rule}\rule{\linewidth}{0.6pt}}

\vspace{8pt}
\begin{tabularx}{\linewidth}{@{}p{2.3cm} L@{}}
Statut & V1 et V2 livrées (partie en direct, spectateur, historique, rejeu, abandon, CI, tournoi suisse). Restent la mise en ligne et le monitoring. \\
Auteur & Raouf Abdallah \\
Document & Référence du projet, du cadrage jusqu'à la mise en ligne. \\
\end{tabularx}

\vspace{8pt}
\tableofcontents
\vspace{4pt}

%---------------------------------------------------------------------
\section{Contexte}

On joue souvent aux échecs entre amis. L'idée de départ est un petit tournoi le même soir, avec un suivi clair pour tout le monde.

Des plateformes existantes couvrent déjà ce besoin de jeu.

\begin{tabularx}{\linewidth}{@{}Y L@{}}
\toprule
Plateforme & Ce qu'elle permet \\
\midrule
Lichess & Tournoi privé gratuit : arène avec mot de passe, ou équipe puis Swiss / arène. \\
Chess.com & Shareable Tournament, lien privé, Swiss ou arène. En général au moins 4 joueurs pour une arène live. \\
Chessly & Apprentissage. Pas d'organisation de tournois. \\
\bottomrule
\end{tabularx}

Ce projet ne cherche pas à remplacer Lichess. Il sert à concevoir un produit temps réel, et à monter en compétence sur la stack la plus présente dans les offres junior en France : Java, Spring Boot, PostgreSQL, Docker, tests et Kafka.

Ce sujet plutôt qu'un CRUD, parce que :

\begin{itemize}
  \item un seul coup est valide à la fois ;
  \item plusieurs écrans restent synchronisés : Blancs, Noirs, spectateur ;
  \item un coup illégal ou hors tour est refusé par le serveur ;
  \item la partie se rejoue depuis un journal d'événements.
\end{itemize}

%---------------------------------------------------------------------
\section{Objectifs}

\subsection{Produit}
\begin{itemize}
  \item Deux amis jouent une partie en direct dans le navigateur.
  \item Un troisième écran suit la partie sans pouvoir jouer.
  \item Un coup illégal est refusé.
  \item La partie se rejoue coup par coup depuis l'historique.
  \item Plus tard : un mini-tournoi entre amis, avec invitation, parties et classement.
\end{itemize}

\subsection{Technique}
Ce que le CV et l'entretien doivent pouvoir défendre, une fois le projet livré.

\begin{itemize}
  \item API Java 21 / Spring Boot 4, source de vérité côté serveur.
  \item PostgreSQL, migrations Flyway, unicité du numéro de coup.
  \item Spring Security et JWT, rôles joueur et spectateur.
  \item Kafka, avec outbox dans la même transaction que le coup.
  \item Second service qui consomme Kafka et pousse le direct.
  \item Docker Compose, JUnit, Mockito, Testcontainers, GitHub Actions.
  \item Front Angular et TypeScript qui affiche l'échiquier, sans décider si un coup est légal.
\end{itemize}

\subsection{Hors objectifs}
\begin{itemize}
  \item Moteur d'analyse : Stockfish, évaluation, flèches.
  \item Classement Elo mondial, anti-triche, matchmaking public.
  \item Kubernetes, Spring Cloud, gRPC, Elasticsearch, Keycloak.
  \item Application mobile native.
  \item Remplacer Chess.com ou Lichess.
\end{itemize}

%---------------------------------------------------------------------
\section{Utilisateurs et rôles}

\begin{tabularx}{\linewidth}{@{}Y L@{}}
\toprule
Rôle & Peut \\
\midrule
Visiteur & Créer un compte, se connecter. \\
Joueur Blanc ou Noir & Rejoindre sa partie, jouer uniquement à son tour. \\
Spectateur & Ouvrir une partie avec son code, suivre les coups, lancer le rejeu. \\
Organisateur (V2) & Créer un tournoi, inviter des amis, voir le classement. \\
\bottomrule
\end{tabularx}

Des comptes de démo sont prévus pour l'entretien : deux joueurs et un spectateur. Les identifiants vivent dans le README, pas dans ce document.

%---------------------------------------------------------------------
\section{Périmètre}

\subsection{V1, partie en direct}
Obligatoire avant le tournoi.

\begin{itemize}
  \item Inscription et connexion JWT.
  \item Créer une partie, obtenir un code, le second joueur rejoint.
  \item Échiquier, coup envoyé au serveur, diffusion aux deux joueurs.
  \item Refus : coup illégal, mauvais tour, spectateur, double envoi du même numéro.
  \item Fin de partie : mat, pat, abandon.
  \item Promotion du pion : le joueur choisit la pièce, le serveur l'applique.
  \item Spectateur en lecture seule, rejeu ordonné.
  \item Lancement local avec Docker Compose, README avec schéma et parcours de démo.
\end{itemize}

\subsection{V2, tournoi entre amis (livrée)}

\begin{itemize}
  \item Créer un tournoi et inviter par code ; l'organisateur est inscrit d'office. 4 à 8 joueurs.
  \item Format retenu : suisse léger, 3 à 5 rondes, ramenées à « joueurs moins un » au lancement si besoin. Pas d'arène (elle demande une pendule) ni de toutes-rondes.
  \item Les parties sont créées par le serveur à chaque ronde, avec leurs deux joueurs et leurs couleurs. La ronde suivante démarre seule quand la dernière partie de la ronde se termine.
  \item Points : victoire 1, nulle 0,5, défaite 0, bye 1 (au plus un par joueur). Départage : Buchholz.
  \item Appariement : ronde 1 au hasard ; ensuite tri par score, sans revanche tant qu'un autre appariement existe (retour arrière), couleurs équilibrées.
  \item Page de tournoi : inscrits, classement, rondes avec lien vers chaque partie, vainqueur, mise à jour en direct.
\end{itemize}

\subsection{Reporté}
\begin{itemize}
  \item Contrôle du temps (blitz, pendule), après une V1 déjà solide.
  \item Chat, notifications e-mail.
  \item Plusieurs parties simultanées au-delà du besoin de démo.
\end{itemize}

%---------------------------------------------------------------------
\section{Parcours de démo}

Trois fenêtres côte à côte.

\begin{enumerate}
  \item Les Blancs jouent un coup.
  \item Les Noirs le voient apparaître.
  \item Le spectateur le voit aussi, et ne peut pas jouer.
  \item Un coup illégal est refusé : l'échiquier ne change pas.
  \item Le rejeu (boutons, flèches du clavier, clic sur un coup) reconstruit la partie coup par coup ; « Revenir à la partie en cours » ramène au direct.
  \item Un joueur abandonne : le résultat s'affiche aux trois écrans.
\end{enumerate}

Si l'on demande pourquoi ce projet :

\quotebox{%
On joue souvent entre amis. Organiser notre tournoi était l'occasion. Lichess et Chess.com savent déjà le faire gratuitement. L'objectif est de concevoir le produit : validation côté serveur, événements, persistance, spectateur et rejeu, sur une stack très présente dans les offres Java en France.%
}

%---------------------------------------------------------------------
\section{Stack}

\subsection{Backend}
\begin{tabularx}{\linewidth}{@{}Y L@{}}
\toprule
Élément & Rôle \\
\midrule
Java 21, Spring Boot 4, Maven & Socle des deux services. \\
Spring Web & API REST. \\
Spring Data JPA, PostgreSQL, Flyway & Persistance et migrations. \\
Spring Security, JWT & Authentification et rôles. \\
Bean Validation & Entrées de l'API. \\
Apache Kafka & Faits de jeu, après l'écriture en base. \\
chesslib & Règles, FEN, coups légaux, mat et pat. \\
\bottomrule
\end{tabularx}

Deux services, pas plus en V1.

\begin{itemize}
  \item \field{game-service} : comptes, parties, coups, persistance, publication.
  \item \field{live-service} : consommation Kafka et push temps réel en SSE. Le rejeu passe par \field{GET /games/\{code\}/moves} sur \field{game-service}.
\end{itemize}

Le serveur n'accepte un coup que s'il figure dans les coups légaux renvoyés par chesslib.

\subsection{Frontend}
\begin{itemize}
  \item Angular, TypeScript, Angular CLI.
  \item chessground pour l'échiquier, dans un composant Angular. Affichage et glisser-déposer uniquement. Il ne valide pas les règles.
  \item Composants, services, routeur et HttpClient pour le REST. SSE (\field{EventSource}) pour le direct.
\end{itemize}

\subsection{Outillage}
Docker, Docker Compose, JUnit 5, Mockito, Testcontainers, GitHub Actions. Éditeur : Cursor, avec les extensions Java et Spring Boot.

OpenAPI et Redis restent optionnels. Redis n'entre que si la présence ou les sessions deviennent un vrai problème.

%---------------------------------------------------------------------
\section{Architecture}

\begin{tabularx}{\linewidth}{@{}Y L@{}}
\toprule
Composant & Rôle \\
\midrule
Navigateur & Angular et chessground. Propose un coup, affiche le FEN reçu. \\
game-service & Décide si le coup est légal. Écrit la partie, le coup et l'outbox. \\
PostgreSQL & Source de vérité : parties, coups, outbox. \\
Kafka & Topic \field{chess.game.events}, clé \field{gameId}. \\
live-service & Consomme les événements et pousse les trois écrans. \\
\bottomrule
\end{tabularx}

\begin{itemize}
  \item Le coup, la mise à jour de la partie et la ligne d'outbox sont dans la même transaction.
  \item Le publisher envoie l'outbox vers Kafka après le commit. Si Kafka est indisponible, le coup reste en base et sera republié.
  \item \field{live-service} ne valide pas les règles. Il diffuse ce que \field{game-service} a déjà accepté.
  \item L'ordre par partie vient de \field{move\_number} en base et de la clé Kafka.
  \item Pas de troisième service tant que ces deux rôles restent clairs.
\end{itemize}

%---------------------------------------------------------------------
\section{Modèle de données}

Migrations uniquement via Flyway. Pas de \field{ddl-auto=update} en dehors d'un essai local jetable.

\subsection{users}
\field{id}, \field{email} unique, \field{password\_hash}, \field{display\_name}, \field{created\_at}.

\subsection{games}
\begin{itemize}
  \item \field{id}, \field{code} unique, court, partageable.
  \item \field{white\_user\_id} et \field{black\_user\_id}, vides jusqu'au join.
  \item \field{status} : \field{WAITING}, \field{IN\_PROGRESS}, \field{FINISHED}.
  \item \field{result} : \field{WHITE\_WIN}, \field{BLACK\_WIN}, \field{DRAW}, ou vide.
  \item \field{end\_reason} : \field{CHECKMATE}, \field{STALEMATE}, \field{RESIGNATION}, ou vide.
  \item \field{current\_fen}, \field{next\_move\_number} (démarre à 1), \field{created\_at}, \field{updated\_at}.
\end{itemize}

\subsection{moves}
\field{id}, \field{game\_id}, \field{move\_number}, \field{player\_color} (\field{WHITE} ou \field{BLACK}), \field{uci} (\field{e2e4}, \field{e7e8q} pour une promotion), \field{san} (\field{e4}, \field{e8=Q}), \field{fen\_after}, \field{created\_at}.

Contrainte unique sur \field{(game\_id, move\_number)}.

\subsection{outbox\_events}
\field{id}, \field{game\_id} (nullable), \field{tournament\_id} (nullable : exactement un des deux est renseigné), \field{event\_type}, \field{payload} JSON, \field{move\_number}, \field{created\_at}, \field{published\_at}. Ce dernier reste vide tant que Kafka n'a pas acquitté.

\subsection{Tournoi (V2)}
\field{tournaments} : \field{id}, \field{code} unique, \field{name}, \field{status} (\field{REGISTRATION}, \field{IN\_PROGRESS}, \field{FINISHED}), \field{created\_by}, \field{rounds\_total}, \field{current\_round}, \field{created\_at}.

\field{tournament\_participants} : \field{id}, \field{tournament\_id}, \field{user\_id}, \field{half\_points} (score en demi-points : victoire 2, nulle 1, pour éviter les décimaux), \field{bye\_round} nullable, \field{joined\_at}. Unique sur \field{(tournament\_id, user\_id)}.

\field{games.tournament\_id} et \field{games.round} sont nullables.

%---------------------------------------------------------------------
\section{API}

Authentification : \field{Authorization: Bearer <jwt>}.

\begin{tabularx}{\linewidth}{@{}p{2.05cm} L L@{}}
\toprule
Méthode & Chemin & Effet \\
\midrule
POST & \field{/auth/register} & Crée un compte. \\
POST & \field{/auth/login} & Renvoie un JWT. \\
POST & \field{/games} & Crée une partie et son code. Blancs par défaut. \\
POST & \field{/games/\{code\}/\allowbreak join} & Le second joueur prend la couleur libre. \\
GET & \field{/games/\{code\}} & FEN, statut, joueurs, dernier coup. \\
POST & \field{/games/\{code\}/\allowbreak moves} & Joue un coup. Corps : \field{uci}. \\
POST & \field{/games/\{code\}/\allowbreak resign} & Abandon du joueur authentifié. \\
GET & \field{/games/\{code\}/\allowbreak moves} & Historique ordonné, pour le rejeu. \\
POST & \field{/tournaments} & Crée un tournoi (\field{name}, \field{rounds} de 3 à 5). L'auteur est organisateur et premier inscrit. \\
GET & \field{/tournaments/\{code\}} & Tournoi, classement (rang, points, Buchholz) et rondes avec leurs parties. \\
POST & \field{/tournaments/\{code\}/\allowbreak join} & Inscription avant le lancement, 8 joueurs maximum. \\
POST & \field{/tournaments/\{code\}/\allowbreak start} & Lance la ronde 1 (organisateur, 4 joueurs minimum). \\
\bottomrule
\end{tabularx}

\vspace{4pt}
Le spectateur s'abonne au flux du \field{live-service}, par \field{gameId} ou par code, sans être assis à l'échiquier.

Codes d'erreur stables, en JSON :

\begin{itemize}
  \item \field{NOT\_YOUR\_TURN}, \field{ILLEGAL\_MOVE}, \field{SPECTATOR\_CANNOT\_MOVE}
  \item \field{GAME\_NOT\_STARTED}, \field{GAME\_ALREADY\_FINISHED}, \field{GAME\_FULL}
  \item \field{TOURNAMENT\_NOT\_FOUND}, \field{TOURNAMENT\_NOT\_OPEN}, \field{TOURNAMENT\_FULL}, \field{ALREADY\_REGISTERED}
  \item \field{NOT\_ORGANIZER}, \field{NOT\_ENOUGH\_PLAYERS}, \field{TOURNAMENT\_ALREADY\_STARTED}
\end{itemize}

Le flux SSE (\field{GET /games/\{code\}/watch}, lecture seule, sans JWT : le code sert de secret de partage) relaie tels quels les messages \field{move.played} et \field{game.ended}. \field{GET /tournaments/\{code\}/watch} fait de même pour un tournoi, sur un canal distinct de celui des parties. À chaque (ré)ouverture, le front recharge l'état par REST pour rattraper ce qu'il a manqué. CORS : origine du front local, puis l'URL de production.

%---------------------------------------------------------------------
\section{Événements Kafka}

Topic \field{chess.game.events}. Clé \field{gameId}, ou \field{t-<tournamentId>} pour un événement de tournoi. Un type par message.

\begin{tabularx}{\linewidth}{@{}Y L@{}}
\toprule
Type & Contenu \\
\midrule
\field{move.played} & gameId, moveNumber, uci, san, fenAfter, playedBy. \\
\field{game.ended} & gameId, moveNumber, result, endReason, finalFen. \\
\field{player.joined} & gameId, color, userId. \\
\field{tournament.updated} & tournamentId, code, status, currentRound. \\
\bottomrule
\end{tabularx}

Le consumer tolère un message déjà vu (rejeu d'offset, republication de l'outbox). Appliquer un coup déjà connu ne duplique pas l'état.

%---------------------------------------------------------------------
\section{Règles côté serveur}

À chaque \field{POST /moves} :

\begin{enumerate}
  \item Charger la partie et la verrouiller pour la transaction.
  \item Vérifier le statut \field{IN\_PROGRESS}, l'identité du joueur, et la couleur attendue.
  \item Charger le FEN dans chesslib.
  \item N'accepter le coup que s'il est dans \field{legalMoves()}. Ne pas se fier à \field{isMoveLegal} seul. Test obligatoire : un pion qui prend droit devant est refusé.
  \item Appliquer le coup, calculer le FEN suivant, détecter mat et pat.
  \item Insérer \field{moves} avec le numéro courant. La contrainte unique protège deux requêtes simultanées.
  \item Mettre à jour \field{games} et insérer \field{outbox\_events}.
  \item Commit. La publication Kafka se fait après, par le publisher.
\end{enumerate}

Promotion : l'UCI contient la pièce, par exemple \field{e7e8q}. Sans pièce de promotion sur la dernière rangée, la réponse est \field{ILLEGAL\_MOVE} et le front redemande le choix.

%---------------------------------------------------------------------
\section{Frontend}

Écrans de la V1 :

\begin{itemize}
  \item Connexion et inscription.
  \item Accueil : créer une partie, rejoindre avec un code, suivre avec un code.
  \item Partie : composant chessground, liste des coups, statut (à vous, à l'adversaire, terminée).
  \item Spectateur : même échiquier, aucun envoi de coup, badge spectateur.
  \item Rejeu : précédent / suivant, à partir de \field{GET /moves}.
  \item Tournoi (V2) : création ou ouverture par code depuis l'accueil ; inscription, lancement par l'organisateur, classement, rondes avec « Jouer ma partie » ou « Suivre », vainqueur. La page de partie propose « Retour au tournoi ».
\end{itemize}

Le message \field{tournament.updated} est un simple signal : il est émis à chaque inscription, lancement, fin de partie de tournoi, changement de ronde et clôture. Le front ne lit pas son contenu, il recharge le tournoi par REST, donc l'ordre des messages n'a pas d'importance.

Le front applique le FEN renvoyé par le serveur. Il peut prévisualiser les cases d'arrivée. Un refus serveur annule le coup visuel.

%---------------------------------------------------------------------
\section{Qualité et tests}

La V1 est présentable quand ces parcours sont couverts. Pas de couverture à 100 \%.

\begin{itemize}
  \item Unitaire : coup illégal, mauvais tour, spectateur, partie déjà finie, promotion.
  \item Concurrence : deux \field{POST /moves} en parallèle. Un seul gagne, un seul \field{move\_number}.
  \item Intégration Testcontainers (PostgreSQL et Kafka) : trois coups produisent trois événements dans l'ordre, et le consumer reconstruit le même FEN que la base.
  \item Le README explique comment lancer les tests.
\end{itemize}

%---------------------------------------------------------------------
\section{Dépôt et construction}

\softbox{%
tournoi-echecs/\\
\quad CAHIER\_DES\_CHARGES.tex\\
\quad README.md\\
\quad docker-compose.yml\\
\quad game-service/\\
\quad live-service/\\
\quad frontend/\\
\quad .github/workflows/ci.yml%
}

L'infrastructure locale est en place : PostgreSQL, Kafka et le topic \field{chess.game.events}. Avancement : \field{game-service} (comptes, JWT, parties, coups, outbox et publisher Kafka), \field{live-service} (consumer Kafka et flux SSE) et le front Angular (jeu en direct, historique SAN, rejeu pas à pas, abandon) existent, et la CI GitHub Actions teste les trois modules. Le rejeu s'appuie sur \field{GET /games/\{code\}/moves}, pas sur le flux. Le tournoi suisse (V2) est livré : inscription par code, rondes et parties créées automatiquement, classement en direct. Restent la mise en ligne et le monitoring. L'ordre de construction ne s'inverse pas.

\begin{enumerate}
  \item Infrastructure : Docker Compose, PostgreSQL, Kafka, topic.
  \item Coup légal en REST et PostgreSQL.
  \item Front Angular qui joue.
  \item Outbox branché sur le Kafka déjà lancé.
  \item \field{live-service}, spectateur, rejeu.
  \item Tests d'intégration et CI.
  \item Tournoi V2.
  \item Mise en ligne.
  \item Monitoring.
\end{enumerate}

Durée indicative : 3 à 4 semaines à mi-temps, ou 2 semaines en continu, une fois Java et Spring en main.

%---------------------------------------------------------------------
\section{Environnement local}

Prérequis : JDK 21, Maven, Node.js LTS, Docker Desktop, Git.

\field{docker compose up} lance aujourd'hui PostgreSQL et Kafka. \field{game-service} et \field{live-service} se lancent à part avec \field{./mvnw spring-boot:run} et rejoindront le même Compose au moment du déploiement. Le front tournera avec \field{ng serve} et un proxy vers les API locales.

Les comptes de démo sont dans le README. Aucun secret réel dans Git.

%---------------------------------------------------------------------
\section{Intégration continue}

GitHub Actions (\field{.github/workflows/ci.yml}), trois jobs à chaque push et pull request :

\begin{itemize}
  \item \field{./mvnw test} sur \field{game-service} et \field{live-service}. Testcontainers utilise le Docker du runner.
  \item \field{npm ci}, \field{ng build}, puis \field{ng test} en Chrome headless sans bac à sable pour le front Angular.
\end{itemize}

À venir : build des images Docker, sans publication tant que le déploiement n'existe pas.

%---------------------------------------------------------------------
\section{Déploiement}

Cible de démo : un VPS (Hetzner, Scaleway ou OVH, environ 5 à 15 euros par mois) avec Docker Compose. Le front statique est sur Vercel ou Cloudflare Pages.

Sur le VPS :

\begin{itemize}
  \item Compose de production : PostgreSQL avec volume, Kafka, les deux services.
  \item Reverse proxy Caddy ou Nginx, HTTPS Let's Encrypt.
  \item PostgreSQL et Kafka non exposés sur Internet.
  \item Secret JWT et mots de passe uniquement en variables d'environnement.
  \item CORS limité à l'URL du front. WebSocket ou SSE autorisés par le proxy.
\end{itemize}

Une fois la V1 stable, GitHub Actions construit les images puis déploie par SSH. Pas de Kubernetes.

Un nom de domaine est optionnel. Une IP ou un sous-domaine suffit pour l'entretien. Un domaine est plus propre pour le CV.

%---------------------------------------------------------------------
\section{Monitoring}

Niveau adapté à une démo, pas à un système d'information bancaire.

\begin{itemize}
  \item Spring Actuator sur les deux services : santé du processus, de PostgreSQL et de Kafka.
  \item Le reverse proxy ne publie pas les endpoints sensibles (environnement, heap dump).
  \item Journaux avec \field{gameId} et \field{moveNumber}. Info sur coup accepté, coup refusé, fin de partie, échec de publication.
  \item Pas de mot de passe, pas de JWT complet dans les logs.
  \item Métriques : coups acceptés et refusés, lag du consumer, taille de l'outbox non publiée.
\end{itemize}

Sur le VPS : \field{docker compose logs}, healthchecks Compose, et une vérification manuelle avant une démo.

Sauvegarde : dump PostgreSQL quotidien tant que des parties comptent. Kafka n'est pas la source de vérité. La base et l'outbox le sont. Si le topic est perdu, on republie l'outbox.

Incident à savoir expliquer : Kafka tombe après le commit. Le coup est en base, \field{published\_at} reste vide, le publisher reprend, le spectateur rattrape. Le joueur a déjà sa réponse HTTP.

%---------------------------------------------------------------------
\section{Sécurité}

Niveau projet personnel.

\begin{itemize}
  \item Mots de passe hashés avec BCrypt.
  \item JWT signé, expiration courte, secret hors de Git.
  \item Un joueur ne joue que sa couleur. Le spectateur n'a pas de route de coup.
  \item Validation des entrées : UCI, longueur du code partie.
  \item Fichier \field{.env} ignoré par Git.
\end{itemize}

Hors V1 : OAuth, double authentification, chiffrement du trafic interne au-delà du réseau Docker.

%---------------------------------------------------------------------
\section{Critères de fin}

La V1 est terminée quand :

\begin{itemize}
  \item la démo des trois écrans marche sur une machine propre, avec le README ;
  \item le coup illégal et le double envoi sont couverts par des tests ;
  \item \field{docker compose up} suffit pour l'API, la base et Kafka ;
  \item le README explique l'architecture en un schéma ;
  \item une URL publique permet de refaire la démo sans le PC du développeur.
\end{itemize}

La V2 est terminée quand quatre comptes jouent un mini-tournoi et qu'un classement s'affiche. C'est le cas : le test d'intégration \field{TournamentApiTests} joue un tournoi de trois rondes à quatre comptes jusqu'au classement final, et le même parcours a été vérifié contre PostgreSQL, Kafka et le flux SSE réels.

%---------------------------------------------------------------------
\section{Formulation CV}

À n'écrire qu'une fois la démo réelle.

\begin{itemize}
  \item Partie d'échecs en direct : API Java 21 / Spring Boot, validation des coups côté serveur, front Angular.
  \item Journal des coups publié sur Kafka via un outbox, ordre garanti par un numéro de coup.
  \item Spectateur et rejeu reconstruits depuis les événements, tests Testcontainers sur les coups illégaux.
\end{itemize}

Java, Spring et Kafka n'entrent dans les compétences maîtrisées qu'à ce moment-là. Le projet reste dans une section Projets, pas dans les expériences professionnelles.

\end{document}
